\section{Cuadernos en general}

Un cuaderno computacional (o \notebook) es\footnote{
  \url{https://docs.jupyter.org/}
}
\blockquote{
  [\ldots]un documento compartible que combina código informático,
  descripciones en prosa, datos, visualizaciones enriquecidas como modelos 3D,
  gráficos, diagramas y figuras, y controles interactivos.
  Un \notebook, junto con un editor [\ldots] proporciona un entorno interactivo rápido
  para construir prototipos y explicar código, explorar y visualizar datos, y compartir
  ideas con otros.
}
Hasta ahora he descrito estos \notebooks~como si fueran una invención nueva,
pero probablemente ya has interactuado con algunos.
Probablemente has visto tutoriales en línea de programación escritos cuadernos de Jupyter.
Si alguna vez has resuelto alguna integral usando Wolfram Mathematica\footnote{
  No confundir con Wolfram Alpha. En vez de un formato de \notebook, Alpha usa una interfaz
  más parecida a la de un motor de búsqueda, procesando una petición a la vez.
} estás familiarizado con la idea general.

Tal vez te preguntes cuál es la diferencia entre escribir código en la forma tradicional
y escribirlo en un \notebook, o por qué se sigue escribiendo código en la forma tradicional
cuando existen los \notebooks.
